\chapterepigraph{No great thing is created suddenly, any more than a bunch of
grapes or a fig.}{Epictetus}{\textit{Discourses}}

\begingroup
\let\clearpage\relax
\let\cleardoublepage\relax
\chapter{The trial wavefunction}
\label{ch:trial_wavefunction}
\endgroup

The trial state is constructed to satisfy fermionic antisymmetry and the electron--electron
cusp analytically, and to keep the learned terms smooth under the Laplacian. Antisymmetry
requires a determinant. The Coulomb divergence admits an analytic cancellation, so it is imposed, not learned. Second derivatives enter the loss, so the features fed to the
network are bounded and smooth and the coordinates are scaled to keep those derivatives
$\mathcal{O}(1)$ across three decades of confinement. The form of the learned part is not
fixed by these constraints; both candidate architectures are constructed here and compared
in Chapter~\ref{ch:kernel}. Unless stated otherwise we work in atomic units.

\section{Preliminaries}
\label{sec:setup}

\paragraph{Hamiltonian and units.}
We study $N$ spin-$\tfrac12$ fermions in two spatial dimensions confined by a harmonic trap
and interacting via Coulomb repulsion,
\begin{equation}
\label{eq:dot-h}
H \;=\; \sum_{i=1}^N \Big( -\tfrac12 \nabla_{\mathbf r_i}^2 + \tfrac12 \,\omega^2 r_i^2 \Big)
\;+\; \sum_{1\le i<j\le N} \frac{1}{|\mathbf r_i-\mathbf r_j|},
\qquad (\hbar=m_e=e=1).
\end{equation}
Spin does not appear in $H$; eigenstates factorize into a spatial function with good $L_z$ and a
spin function with good $(S,S_z)$.

\paragraph{Coordinates and scaling.}
We collect positions as $R=(\mathbf r_1,\ldots,\mathbf r_N)\in\mathbb R^{2N}$ and use trap scaling
$\tilde{\mathbf r}_i=\sqrt{\omega}\,\mathbf r_i$, $\tilde{R}=\sqrt{\omega}\,R$. This aligns typical
length scales to $\mathcal O(1)$, improves conditioning of second derivatives, and makes the
noninteracting ground state $\omega$-independent. When useful, we also use pair distances
$r_{ij}=|\mathbf r_i-\mathbf r_j|$ and their scaled versions $\tilde r_{ij}=|\tilde{\mathbf r}_i-\tilde{\mathbf r}_j|$.

\paragraph{Sectors and focus.}
We primarily report \emph{closed-shell} cases (unique spin singlet at the noninteracting level).
Open shells introduce a degenerate manifold and additional choices for the reference; we keep this
to brief remarks only where needed.

\paragraph{Spin–orbitals and basis.}
One-particle spatial orbitals are the Cartesian two-dimensional harmonic-oscillator
eigenstates $\{\phi_{n_x n_y}(\mathbf r)\}$ with $E_{n_x n_y}=\omega(n_x+n_y+1)$,
$n_x,n_y\in\mathbb N_0$ (the basis used in the implementation; for the closed shells studied
it spans the same filled shells as any complete basis of those shells). Spin–orbitals
are $\varepsilon_\alpha(x)=\phi_{n_x n_y}(\mathbf r)\chi_\sigma(s)$, $\sigma\in\{\uparrow,\downarrow\}$.

% ==========================================================

\section{Overview of the ansatz}
\label{sec:ansatz}
Our trial state combines exact antisymmetry, an analytic short-range cusp, and a compact
neural residual, while optionally evaluating the Slater determinant on backflowed coordinates:
\begin{align}
  \Psi_{\theta,\beta}(R)
  &= \mathrm{SD}\big(\tilde{R}'\big)\;
     \exp\!\Big(
       \sum_{i<j} u_{\sigma_i \sigma_j}(r_{ij})
       + W_\theta(\tilde{R})
     \Big),\label{eq:ansatz-overall} \\
  \tilde{R}' &= \tilde{R} + \Delta_\beta(\tilde{R}).\label{eq:bf-coords}
\end{align}
Here $\mathrm{SD}$ enforces antisymmetry; $u$ hard-wires the Kato cusp; $W_\theta$ (the NQS correlator) learns smooth
medium/long-range structure; and $\Delta_\beta$ (backflow) adjusts the nodes by smoothly warping the coordinates used
inside the determinants; its per-particle displacements $\Delta_{\beta,i}$ are written
$\Delta\mathbf{x}_i$ in the results chapters. We deliberately feed the NQS with $\tilde{R}$ (not $\tilde{R}'$) to avoid self-reinforcing feedback
and keep local-energy variance low.

\paragraph{Design principles.}
(i) \emph{Separate guaranteed physics from learned details:} antisymmetry and the cusp are hard-wired; the network learns the rest.
(ii) \emph{Condition the geometry:} work in trap units and use soft-core pair features so that first/second derivatives are bounded.
(iii) \emph{Keep the learned object small and smooth:} once the cusp is explicit, a compact NQS suffices across $(N,\omega)$.

\section{Slater reference}
\label{subsec:slater}
For closed shells we use a restricted product of spin-block determinants of the lowest $N/2$ harmonic oscillator orbitals (trap units):
\begin{equation}
  \mathrm{SD}(\tilde{R})
  =
  \det\!\big[\phi_a(\tilde{\mathbf r}_i)\big]_{i\in\uparrow,\;a\le N/2}\;
  \det\!\big[\phi_a(\tilde{\mathbf r}_j)\big]_{j\in\downarrow,\;a\le N/2}.
  \label{eq:slater}
\end{equation}
The correlator reweights amplitudes \emph{within} nodal pockets; without backflow, the nodes are those of~\eqref{eq:slater}.

\section{Neural correlator $W_\theta$}
\label{subsec:nqs}
As with the backflow, the correlator comes in two forms that are compared directly in
\S\ref{sec:kernel-q1}: a \emph{separable} network that encodes particle- and
pair-features independently and pools them, and a copresheaf \emph{message-passing}
(CTNN) network. We describe the separable network first, since it is the baseline and supplies the safe features and analytic cusp shared by both, and then the message-passing one.

\subsection{Separable correlator (DeepSets baseline)}
\label{subsec:separable-correlator}
The separable NQS aggregates \emph{particlewise} and \emph{pairwise} embeddings with permutation-invariant pooling and appends two global
statistics, all in $\tilde{R}$. A core safety requirement is that pair channels satisfy $d s_k/d\tilde r \to 0$ as $\tilde r\to 0$
so that $W_\theta$ never injects spurious $1/\tilde r$ terms into $\nabla\log\Psi$ or $\Delta\log\Psi$.

\subsubsection{Particle branch (DeepSets).}
This part of the network is typically referred to as Deep Sets~\cite{Zaheer2017-DeepSets}. The permutation invariance of quantum many-body mechanics is naturally captured by this architecture. 
A neural network $\phi$ creates a latent embedding for each particle individually. All embeddings are pooled together to form a global representation, which is used later in the readout. Formally, we have
\[
  \mathbf h_i^\phi=\phi(\tilde{\mathbf x}_i)\in\mathbb R^{d_L},\qquad
  \overline{\mathbf h}^{\phi}=N^{-1}\sum_i \mathbf h_i^\phi\in\mathbb R^{d_L}.
\]

\subsubsection{Pair branch (soft-core, even channels).}
Although the particle branch could theoretically produce pair correlations, in practice it struggles to represent sharp near-field structure. We therefore add a dedicated pair branch that builds features from interparticle distances.
This network processes all pairs $(i,j), i\ne j$ individually and pools the resulting latent embeddings. 
The problem is that raw distances $\tilde r_{ij}$ have unbounded derivatives at coalescence, which destabilizes Laplacian evaluations. We therefore build \emph{safe} pair features with bounded first and second derivatives as $\tilde r\to 0$.
With $\tilde r^{\rm soft}=\sqrt{\tilde r^2+\varepsilon^2}$ ($\varepsilon\!\in[0.15,0.30]$ in trap units), we build six “safe” features
\begin{align}
  s_1 &= \log\!\big(1+(\tilde r^{\rm soft}/\varepsilon)^2\big),\quad
  s_2 = \tfrac{\tilde r^{2}}{\tilde r^{2}+\varepsilon^2},\quad
  s_3 = (\tilde r^{\rm soft}/\varepsilon)^{2}\exp[-(\tilde r^{\rm soft}/\varepsilon)^{2}], \notag\\
  \mathrm{rbf}_g &= \exp(-g\,s_1),\quad g\in\{0.25,1,4\},\qquad
  \mathbf u_{ij}=[s_1,s_2,s_3,\mathrm{rbf}_{0.25},\mathrm{rbf}_{1},\mathrm{rbf}_4]\in\mathbb R^{6}.
  \label{eq:safe-feats}
\end{align}
An MLP $\psi$ maps $\mathbf u_{ij}\mapsto\mathbf h^{\psi}_{ij}\in\mathbb R^{d_L}$. A short-range gate
\[
  \chi(\tilde r)=\frac{\tilde r^{2}}{\tilde r^{2}+r_g^{2}},\quad r_g\approx 0.3,\quad \chi(0)=\chi'(0)=0,
\]
suppresses learned responses \emph{at} coalescence; downstream we use
$\chi\,\mathbf h^\psi_{ij}$.

\subsection{Analytic short-range cusp (shared)}
\label{subsec:cusp-implementation}
The analytic cusp below is shared by both correlator variants (and is independent of the
backflow). Because the safe features are necessary in the pair branch, the NQS cannot represent the Kato cusp condition exactly.
We therefore add an analytic cusp term to the exponent in~\eqref{eq:ansatz-overall},
written in \emph{physical} coordinates so that its slope at contact is manifestly
correct:
\begin{equation}
  u_{\sigma\sigma'}(r)
  =
  \gamma_{\sigma\sigma'}\, r\, e^{-r/a_{\rm ho}},
  \qquad a_{\rm ho}=1/\sqrt{\omega},
  \qquad
  \gamma_{\uparrow\downarrow}=\tfrac{1}{d-1},\quad
  \gamma_{\uparrow\uparrow}=\tfrac{1}{d+1}.
  \label{eq:cusp}
\end{equation}
Because the linear prefactor is the physical separation $r$, the slope at contact is
$\partial_r u|_{0}=\gamma_{\sigma\sigma'}$, exactly the Kato condition in general $d$;
for $d=2$ this gives $\gamma_{\uparrow\downarrow}=1$ and $\gamma_{\uparrow\uparrow}=\tfrac13$,
the standard two-dimensional quantum-dot values, derived---including the parallel-spin
case, where antisymmetry already forces a linear node---in
Appendix~\ref{app:coulomb:cusp}. Only the exponential taper carries the trap length
$a_{\rm ho}$, suppressing long-range interference so the NQS need not unlearn cusp
structure at medium and long range; in trap units ($\tilde r=\sqrt{\omega}\,r$) the same
term reads $u=(\gamma_{\sigma\sigma'}/\sqrt{\omega})\,\tilde r\,e^{-\tilde r}$.


\subsection{Separable readout}
With contributions from the particle and pair branches, we form a global representation for the final readout. In addition to the pooled embeddings $\overline{\mathbf h}^\phi$ and $\overline{\mathbf h}^\psi$, we also include two global statistics that help the network gauge overall system size and density.
We average pairs by default or apply degree-normalised radial attention
\[
  w_{ij}=\exp[-(\tilde r_{ij}/r_c)^p],\quad
  \hat w_{ij}=w_{ij}\Big/\sum_{k\ne i}w_{ik},\quad
  (r_c,p)=(0.4,6),
\]
then form $g_i=\sum_{j\ne i}\hat w_{ij}\,\mathbf h^\psi_{ij}$ and $\overline{\mathbf h}^\psi=N^{-1}\sum_i g_i$.
Two system-level scalars in trap units,
\[
  \mathbf g(\tilde{R})=
  \Big[\tfrac{1}{Nd}\sum_i\|\tilde{\mathbf x}_i\|^2,\;
       \tfrac{2}{N(N-1)}\sum_{i<j}s_1(\tilde r_{ij})\Big]\in\mathbb R^2,
\]
are concatenated with $\overline{\mathbf h}^{\phi}$ and $\overline{\mathbf h}^{\psi}$ and mapped to a scalar by a tiny MLP $\rho$:
\begin{equation}
  W_\theta(\tilde{R})
  =
  \rho\!\Big(\overline{\mathbf h}^{\phi}\ \|\ \overline{\mathbf h}^{\psi}\ \|\ \mathbf g(\tilde{R})\Big).
  \label{eq:nqs-readout}
\end{equation}

\subsection{Message-passing correlator (CTNN)}
\label{subsec:ctnn-correlator}
The message-passing correlator keeps the same node/edge
embeddings, safe pair features, and analytic cusp, but replaces the independent pooling of
the separable network with the copresheaf transport of \S\ref{sec:ctnn-theory}: at each
round the endpoint nodes are transported into an edge update and the updated edges are
transported back and aggregated onto the nodes,
\begin{equation}
  \mathbf h_{ij} \mathrel{+}= F_e\!\big(\mathbf h_{ij}\,\|\,\rho_{v\to e}\mathbf h_i\,\|\,\rho_{v\to e}\mathbf h_j\big),
  \qquad
  \mathbf h_i \mathrel{+}= F_v\!\Big(\mathbf h_i\,\|\,\operatorname*{Agg}_{j\neq i} w_{ij}^{\rm spin}\,\rho_{e\to v}\mathbf h_{ij}\Big),
  \label{eq:ctnn-corr-step}
\end{equation}
with residual updates and per-round transport maps. Unlike the single-round backflow, the
correlator stacks these rounds in a multiscale ``V-cycle'': $n_{\downarrow}$ down-pass
rounds, a linear coarsening to a bottleneck width, then $n_{\uparrow}$ up-pass rounds with
skip connections back to the down-pass features (defaults $n_{\downarrow}=n_{\uparrow}=2$).
A permutation-invariant readout over the final node and edge statistics returns the scalar
$W_\theta(\tilde R)$, exactly as in~\eqref{eq:nqs-readout}. Holding the Slater reference,
cusp, and training fixed and swapping only the separable readout for this V-cycle isolates
the relational channel analysed in \S\ref{sec:kernel-q1}.

\section{Backflow}
\label{subsec:backflow-method}

Backflow deforms the coordinates \emph{used only inside the Slater reference} so that the nodal surface can adapt to interparticle correlations:
\begin{equation}
  \tilde{R}' \;=\; \tilde{R} + \Delta_\beta(\tilde{R}) ,
  \qquad
  \Psi_{\theta,\beta}(R) \;=\; \mathrm{SD}(\tilde R')\;
  \exp\!\Big(\sum_{i<j} u_{\sigma_i\sigma_j}(r_{ij}) + W_\theta(\tilde R)\Big).
  \label{eq:bf-def}
\end{equation}
The correlator $W_\theta$ and cusp $u$ are evaluated on $\tilde R$ to avoid feedback loops in the local energy; only the determinant ``sees'' $\tilde R'$.
This choice cleanly separates (i) amplitude reweighting within nodal pockets (handled by $u+W_\theta$) from (ii) nodal \emph{geometry} (handled by backflow).

\paragraph{Design goals.}
(i) \emph{Stability near coalescence:} messages are built from mollified radii and
short-range gates so that $\partial \Delta_\beta/\partial \tilde r$ stays bounded;
(ii) \emph{Permutation symmetry:} permutation equivariance holds by construction (shared
maps and symmetric aggregation);
(iii) \emph{Close to identity at initialisation:} a zero-initialised last layer and a
positive learnable scale keep $\Delta_\beta\!\approx\!0$ at the start.
Rotational symmetry, by contrast, is only \emph{approximate}: the displacement head reads
out a $d$-vector from node features that include the raw coordinates, so the map is not
rotation-equivariant by construction, and it is encouraged instead by the
geometry-based edge features and by rotationally-augmented sampling. A global translation
is not a symmetry of the trap in any case; the zero-mean projection~\eqref{eq:bf-com} is a
regulariser that removes a spurious centre-of-mass drift, not an imposed physical symmetry.

We study \emph{two} backflow architectures. Both map the configuration to a per-particle
displacement field and share the output bounding, spin masking, and centre-of-mass
projection below; they differ only in how a particle's displacement is made to depend on
the others. The conventional network is the baseline of the architecture comparison in
Chapter~\ref{ch:kernel}; the copresheaf \textsc{CTNN} is the production ansatz.

\subsection{Conventional backflow (single-round messages)}
The baseline uses a single round of pairwise messages. For
each ordered pair $(i,j)$ with $i\neq j$,
\begin{equation}
  \mathbf m_{ij}
  = \phi_\beta\!\big(
      \tilde{\mathbf x}_i,\ \tilde{\mathbf x}_j,\ \mathbf r_{ij},\
      \tilde r_{ij}^{\rm soft},\ (\tilde r_{ij}^{\rm soft})^2
    \big)\cdot w_{ij}^{\rm spin},
  \qquad \tilde r_{ij}^{\rm soft}=\sqrt{\tilde r_{ij}^2+\varepsilon_{\rm bf}^2},
  \label{eq:bf-message}
\end{equation}
with $\phi_\beta$ a small MLP and $w_{ij}^{\rm spin}\in\{0,1\}$ a spin mask (same-spin only
or all pairs). Messages are aggregated to a node state and a node network returns the
displacement,
\begin{equation}
  \mathbf m_i = \operatorname*{Agg}_{j\neq i}\mathbf m_{ij},
  \qquad
  \delta\tilde{\mathbf x}_i = \psi_\beta\!\big(\tilde{\mathbf x}_i \,\|\, \mathbf m_i\big),
  \qquad \mathrm{Agg}\in\{\mathrm{sum},\mathrm{mean},\mathrm{max}\}.
\end{equation}
There is no persistent edge state: each message is recomputed from the raw geometry and
discarded after aggregation.

\subsection{CTNN backflow (copresheaf)}
\label{subsec:ctnn-backflow}
The message-passing architecture instead maintains
explicit \emph{node and edge} features and transports information between them through
learned linear maps---the copresheaf-style step of \S\ref{sec:ctnn-theory}. Node features
are embedded from the position and a spin scalar, edge features from the relative geometry,
\begin{equation}
  \mathbf h_i^{(0)} = W_v\,[\tilde{\mathbf x}_i,\, s_i],
  \qquad
  \mathbf h_{ij}^{(0)} = \mathrm{MLP}_e\!\big(\mathbf r_{ij},\, |\mathbf r_{ij}|,\, |\mathbf r_{ij}|^2\big).
\end{equation}
A node-to-edge transport $\rho_{v\to e}$ feeds the endpoint nodes into an edge update; an
edge-to-node transport $\rho_{e\to v}$ returns the updated edges to the nodes, which are
refined by a residual node update; a linear head reads out the displacement:
\begin{align}
  \mathbf h_{ij}' &= \mathrm{MLP}_e'\!\big(\mathbf h_{ij}^{(0)}\,\|\,\rho_{v\to e}\mathbf h_i^{(0)}\,\|\,\rho_{v\to e}\mathbf h_j^{(0)}\big),
  \label{eq:ctnn-edge}\\
  \mathbf m_i &= \operatorname*{Agg}_{j\neq i} w_{ij}^{\rm spin}\,\rho_{e\to v}\big(\mathbf h_{ij}'\big),
  \label{eq:ctnn-agg}\\
  \mathbf h_i &= \mathbf h_i^{(0)} + \mathrm{MLP}_v\!\big(\mathbf h_i^{(0)}\,\|\,\mathbf m_i\big),
  \qquad
  \delta\tilde{\mathbf x}_i = W_{\!\Delta}\,\mathbf h_i.
  \label{eq:ctnn-node}
\end{align}
The learned transport maps $\rho_{v\to e},\rho_{e\to v}$ are the channel whose ablation
collapses the backflow at the Wigner crystal (Chapter~\ref{ch:kernel}). The backflow uses a
single copresheaf round; the message-passing \emph{correlator} of the architecture
comparison (\texttt{CTNNJastrowVCycle}) iterates the same style of update over a down/up
``V-cycle''.

\subsection{Constraints, symmetries and scaling}
\label{sec:bf-init}
Both variants take inputs in trap units ($\tilde{\mathbf x}=\sqrt{\omega}\,\mathbf x$), bound
the displacement, and start at the Slater nodes,
\begin{equation}
  \Delta_\beta(\tilde R)\;=\;\tanh\!\big(\delta\tilde X\big)\cdot s_\beta,
  \qquad s_\beta=\mathrm{softplus}(s_\beta^{\rm raw})>0,
  \label{eq:bf-bound-scale}
\end{equation}
with the last linear layer zero-initialised so that $\Delta_\beta\!\approx\!0$ initially and
$s_\beta$ initialised in $[0.02,0.10]$. This parameterisation stabilises the SR and energy
tails and protects the Laplacian.

Near coalescence raw $1/\tilde r$ factors make $\nabla\!\cdot\!\Delta_\beta$ spiky. We
mollify $\tilde r_{ij}$ as in \eqref{eq:bf-message} and optionally gate the message,
\begin{equation}
  \chi_{\rm bf}(\tilde r_{ij})=\frac{\tilde r_{ij}^2}{\tilde r_{ij}^2+r_g^2},
  \qquad \chi_{\rm bf}(0)=\chi'_{\rm bf}(0)=0,\quad r_g\approx 0.3,
  \label{eq:bf-gate}
\end{equation}
or use degree-normalised radial attention $\hat w_{ij}=w_{ij}/\sum_{k\neq i}w_{ik}$ with
$w_{ij}=\exp[-(\tilde r_{ij}/r_c)^p]$. Either keeps the learned response smooth at contact and
limits the backflow's contribution to the local-energy variance. The mollification scale
$\varepsilon_{\rm bf}=10^{-6}$ is fixed, not tuned, and is a numerical guard rather
than a modelling choice: $\sqrt{\tilde r^2+\varepsilon_{\rm bf}^2}$ differs from $\tilde r$ by
less than $\varepsilon_{\rm bf}$ everywhere, so it changes no value the network can resolve
and exists only to keep $\nabla\tilde r$ finite at exact coincidence. The physical short-range
scale is instead the pair-feature softening $\varepsilon\in[0.15,0.30]$ of
\eqref{eq:safe-feats}, four to five orders of magnitude larger.

Permutation symmetry is automatic. To remove a spurious centre-of-mass drift the flow is
projected to zero mean,
\begin{equation}
  \Delta_\beta \;\leftarrow\; \Delta_\beta \;-\; \frac{1}{N}\sum_{i=1}^N \Delta_{\beta,i},
  \label{eq:bf-com}
\end{equation}
taken over particles within each configuration, not across the batch. Messages are
built from relative vectors $\mathbf r_{ij}$, but the map is not translation-invariant or
rotation-equivariant by construction, since the displacement head reads a $d$-vector from
absolute coordinates. The two variants apply the projection at different points: the
copresheaf network applies \eqref{eq:bf-com} inside its forward pass, while the conventional
network returns the raw displacement and the caller projects when $\tilde R'$ is formed. The
shifted coordinates are identical either way, but a diagnostic evaluated on a network's
\emph{returned} field sees a centre-of-mass component in one case and not the other, which is
why the two rank columns of \S\ref{sec:kernel-q1b} sit under ceilings differing by two.

Because only the determinant is evaluated at $\tilde R'$, backflow moves the \emph{nodes}
while the correlator exponent stays on $\tilde R$, and no Jacobian term arises. This
Feynman--Cohen construction~\cite{FeynmanCohen1956-Backflow} targets fixed-node error
directly; its modern quantum-Monte-Carlo form is due to Holzmann and
Ceperley~\cite{HolzmannCeperley2003-Backflow}. Compute and memory scale as
$O(BN^2(d+H))$ for batch size $B$ and hidden width $H$, which is what limits $N$ in practice.


\section{Architecture sizes and hyperparameters}
\label{sec:hyperparameters}
Table~\ref{tab:hyperparams} collects the representative settings. The two comparisons in
Chapter~\ref{ch:kernel} stand on different footings in this respect.

The \emph{correlator} comparison is capacity-controlled. The two variants are held at
comparable parameter counts, so a difference in outcome cannot be attributed to one simply
being larger, and where \S\ref{sec:kernel-q1a} departs from matched capacity it does so in
the \emph{separable} network's favour, scaling it from $20$k to $164$k parameters against a
fixed $80$k message-passing network without closing the gap.

The \emph{backflow} comparison is not. At equal width the copresheaf network is intrinsically
the heavier object---it carries a persistent edge state where the conventional network carries
only a per-particle aggregate---so at the shared production setting (message width $128$,
node MLP $128\times3$) it holds $182$k parameters against $51$k, a factor of $3.6$. No
matched backflow pair was trained. What separates architecture from capacity there is
therefore not a matched control but the conventional network's own behaviour across
$\omega$, which \S\ref{sec:kernel-q1b} sets out where the measurement is reported. All settings below are those of the accompanying code
the accompanying code.

\begin{table}[htbp]\centering
\caption[Architecture and training hyperparameters]{The three sets of models trained in this
thesis. They exist for different purposes and are never compared with one another. Widths are
hidden dimensions; ``rounds'' are message-passing steps ($n_\downarrow{+}n_\uparrow$ for the
V-cycle). Within the architecture study the correlator variants are not capacity-matched,
which is why \S\ref{sec:kernel-q1a} also reports the matched $89.4$k separable network and
the full $20$k--$164$k ladder; the backflow variants are not matched either, and cannot be at
equal width, for the reason given above. Settings follow the accompanying code
(\texttt{scripts/run\_depth\_analysis.py}, \texttt{scripts/source\_state\_check.py}).}
\label{tab:hyperparams}
\setlength{\tabcolsep}{5pt}
\renewcommand{\arraystretch}{1.25}
\begin{tabular}{@{}p{0.165\linewidth} p{0.335\linewidth} p{0.30\linewidth} r@{}}
\hline
Programme & Correlator & Backflow & Total \\
\hline
Physics production (\S\ref{sec:energy-overview})
  & separable; width $128$, $2$ layers; $54$k
  & conventional; msg $128$, node $128\times3$; $51$k & $105$k \\
  & \emph{(as above)} & message-passing; same widths; $182$k & $236$k \\
\hline
Architecture study (\S\ref{sec:kernel-q1})
  & separable; pair-MLP $64\times4$, out $32$; $47.7$k
  & \multirow{2}{*}{\emph{as production; not matched}} & --- \\
  & message-passing; node/edge $32$, bottleneck $16$, $2{+}2$ rounds; $79.8$k & & --- \\
\hline
Collocation (\S\ref{sec:kernel-mcmcfree})
  & message-passing V-cycle; $25$k & none & $25$k \\
  & \emph{(as above)} & copresheaf; $h_{\rm bf}{=}128$; $+24$k & $49$k \\
\hline
\end{tabular}

\medskip
\small\noindent Readout for every correlator: hidden $64$, $2$--$3$ layers, SiLU.
Optimiser: Adam at $\mathrm{lr}=5\times10^{-4}$ (backflow) and $5\times10^{-5}$ (Jastrow);
CG-SR with Tikhonov damping $\lambda$, a trust-region cap and CG restarts. Collocation batch
$n_{\rm coll}\sim3\times10^{3}$ with oversampling $m=8$--$16$; the collocation families are
described further in \S\ref{subsec:colloc-families}.
\end{table}


\paragraph{Implementation correspondence.} For replication: the separable correlator is
\texttt{PINN} and the message-passing one \texttt{CTNNJastrowVCycle}; the conventional
backflow is \texttt{BackflowNet}, whose $\phi_\beta$ and $\psi_\beta$ are its message and
node MLPs, and the copresheaf backflow is \texttt{CTNNBackflowNet}, whose embeddings,
transport maps and update MLPs realise \eqref{eq:ctnn-edge}--\eqref{eq:ctnn-node}.
Aggregation is selectable (\texttt{sum|mean|max}); the bounded output of
\eqref{eq:bf-bound-scale} is \texttt{out\_bound="tanh"} with
$s_\beta=\mathrm{softplus}(\texttt{bf\_scale\_raw})$. The comparison tiers are produced by
\texttt{scripts/run\_depth\_analysis.py} and the production widths by
\texttt{scripts/source\_state\_check.py}. Schedule and solver options carry the names
\texttt{alpha\_start}, \texttt{alpha\_end}, \texttt{alpha\_decay\_frac},
\texttt{max\_param\_step}, \texttt{damping}, \texttt{micro\_batch} and \texttt{grad\_clip}.

\chapter{Optimisation}
\label{ch:optimisation}
\label{ch:optimization}

\section{Primary objective and two–stage training}
\label{sec:opt-overview}

Given a parameterised wavefunction $\Psi_\theta$, the variational ground–state energy is
\begin{equation}
E(\theta)
=\frac{\langle\Psi_\theta|H|\Psi_\theta\rangle}{\langle\Psi_\theta|\Psi_\theta\rangle}
=\mathbb{E}_{\tilde{R}\sim \pi_\theta}\!\big[E_L(\tilde{R})\big],
\qquad
\pi_\theta(\tilde{R})=\frac{|\Psi_\theta(\tilde{R})|^2}{\int |\Psi_\theta|^2}.
\end{equation}
The local energy is evaluated in \emph{physical} coordinates, consistently with the
Hamiltonian~\eqref{eq:dot-h},
\begin{equation}
\label{eq:local-energy}
E_L(R)
= -\tfrac12\sum_{i=1}^N\!\Big[\Delta_{\mathbf r_i}\ln\Psi_\theta + \|\nabla_{\mathbf r_i}\ln\Psi_\theta\|^2\Big]
+\tfrac{1}{2}\,\omega^2\sum_{i=1}^N r_i^2
+\sum_{i<j}\frac{1}{r_{ij}},
\end{equation}
and the zero–variance principle holds: $\mathrm{Var}_{\pi_\theta}[E_L]=0$ iff $\Psi_\theta$ is an eigenstate.
The network \emph{features} are computed in trap units $\tilde{\mathbf r}=\sqrt{\omega}\,\mathbf r$
for conditioning, but the local energy itself is formed in physical units, so no
$\omega$ factors are absorbed into the kinetic, trap, or Coulomb terms. (In trap units
the same operator reads $E_L=\omega\big[-\tfrac12\sum_i(\Delta_{\tilde{\mathbf r}_i}\ln\Psi_\theta+\|\nabla_{\tilde{\mathbf r}_i}\ln\Psi_\theta\|^2)+\tfrac12\sum_i\tilde r_i^2\big]+\sum_{i<j}\sqrt{\omega}/\tilde r_{ij}$;
the kinetic and trap terms scale by $\omega$ and the Coulomb term by $\sqrt{\omega}$.)

Our optimisation follows two stages that mirror the implementation:
\begin{enumerate}
  \item \textbf{Residual–based pretraining (stage I).} We fit the residual of $E_L$ to a moving scalar target $E_{\rm eff}$ using the loss
  \(
    \mathcal{L}_{\rm res}(\theta)=\mathbb{E}\!\big[(E_L(\tilde R)-E_{\rm eff})^2\big].
  \)
  We begin with a \emph{self–target} (stabilizer) $E_{\rm eff}=\mu$ where $\mu=\mathbb{E}[E_L]$ over the current batch. This minimises the \emph{sample} variance of $E_L$ under the current collocation measure---which is not the same object as $\Var_{|\Psi|^2}(E_L)$ unless the two measures coincide---and prevents the network from chasing a potentially inaccurate external target before it can represent low–variance pockets. We then anneal toward a trusted reference $E_{\rm DMC}$ via
  \begin{equation}
    E_{\rm eff}(\alpha) \;=\; \alpha\,E_{\rm DMC} + (1-\alpha)\,\mu,
    \qquad \alpha \in [0,1],
  \end{equation}
  with a smooth cosine schedule $\alpha=\alpha(t)$ that ramps from $\alpha_{\rm start}$ to $\alpha_{\rm end}$ over a prescribed fraction of epochs. Setting the \texttt{objective} to \texttt{"residual"} uses $E_{\rm eff}=\mu$; \texttt{"energy"} uses $E_{\rm eff}=E_{\rm DMC}$; and \texttt{"energy\_var"} uses the annealed blend above.
  \item \textbf{Stochastic Reconfiguration (stage II).} After residual pretraining, we refine with SR (natural–gradient VMC). Writing $\bar O(\tilde R)=\partial_\theta \log\Psi_\theta(\tilde R)-\mathbb{E}_\pi[\partial_\theta\log\Psi_\theta]$ for the centred score,
  \begin{equation}
    S \;=\; \mathbb{E}_\pi\!\big[\bar O^\top \bar O\big],\qquad
    g \;=\; 2\,\mathbb{E}_\pi\!\big[(E_L-\mathbb{E}_\pi[E_L])\,\bar O\big],
  \end{equation}
  we solve the damped normal equations
  \begin{equation}
    (S+\lambda I)\,\Delta\theta \;=\; -\,g,
  \end{equation}
  using (preconditioned) conjugate gradients with restarts, a trust–region scale on the final step, and batchwise centering/whitening of $O$.
\end{enumerate}
This pipeline exploits the zero–variance principle for stable shaping of $W_\theta$ (and backflow) before applying the geometry–aware SR update.

\section{Configuration–space sampling}
\label{sec:sampling}

\paragraph{Guiding principle (permutation invariance).}
In a many–body fermionic system, \emph{which} particle indices are close is irrelevant; what matters is \emph{how many} close pairs and at which length scales. Consequently, we can design samplers that (i) emphasize \emph{configurational motifs} (clusters, shells, dimers, long tails) and not specific labels, and (ii) deliberately permute particles and randomly rotate configurations to explore the space efficiently without biasing towards any ordering.

\subsection{Stratified capped–simplex mixture}
We draw collocation batches $X\in\mathbb{R}^{B\times N\times d}$ from a five–component mixture,
\[
\text{\small(centre, tails, mixed, shells, dimers)},
\]
with nonnegative weights $w\in\Delta^4$ projected to a \emph{capped simplex} $\{\,w:\, \sum_k w_k=1,\ 0\le w_k\le 0.3\,\}$ to enforce coverage.
Each component proposes $x$ as follows (all in physical units; the code internally scales by $a_{\rm ho}=1/\sqrt{\omega}$):
\begin{itemize}
  \item \textbf{Centre:} narrow isotropic Gaussian around the trap centre (samples near the density core).
  \item \textbf{Tails:} wide Gaussian (probes low–density outskirts and large–$r$ behaviour).
  \item \textbf{Mixed:} hybrid batches mixing narrow and moderate spreads per configuration (tests robustness to inhomogeneous spreads).
  \item \textbf{Shells:} points placed on $K$ hyperspherical shells with trap–scaled radii $\{r_k\}$ and occupancy $\{q_k\}$, with small radial/tangential jitter (captures ring/shell order at weak traps).
  \item \textbf{Dimers:} we induce a few close pairs per configuration by sampling short interparticle offsets with random directions and log–uniform distances (enforces near–coalescence coverage without singularities).
\end{itemize}
We then apply a random in–plane rotation (for $d=2$), and a random permutation of particle indices; both preserve the target expectations by symmetry.

\paragraph{Adaptive mixture and shell occupancy.}
Per epoch we compute simple difficulty scores and update $w$ and $q$ by exponentiated–gradient (EG) steps with momentum, temperature, and an exploration term; $w$ is then projected back to the capped simplex. Mixture difficulty uses the residual $(E_L-E_{\rm eff})^2$ per component; shell difficulty uses a proxy
\[
\underbrace{V_i}_{\text{harm.+Coulomb per particle}}
\;+\; \underbrace{\gamma_{g^2}\,\|\nabla_{\mathbf x_i}\log\Psi\|^2}_{\text{stiff gradients}}
\;+\; \underbrace{\gamma_{\rm cusp}\,r_{\min}^{-1}}_{\text{coalescence pressure}},
\]
aggregated by nearest shell index. This shifts probability mass toward underfit regions \emph{without} collapsing coverage thanks to the caps.

\paragraph{Hard–example injection.}
We maintain a small buffer of previously challenging configurations (based on residuals) and stochastically inject them, with mild multiplicative and additive jitter, into the next batch. This provides a curriculum–like pressure while avoiding overfitting to outliers.

\paragraph{Robustness tweaks.}
We trim extreme local–energy outliers by quantiles (default $3\%$) \emph{before} forming losses/updates; we also micro–batch the evaluation to keep memory bounded, and apply gradient clipping to $f_\theta$ (and backflow) parameters.

\section{Residual–based training (stage I)}
\label{sec:residual-stage}

For a batch $\{\tilde R_b\}_{b=1}^B$ we compute $E_L(\tilde R_b)$ and the batch mean $\mu=\frac{1}{B}\sum_b E_L(\tilde R_b)$.
The residual loss is
\begin{equation}
\mathcal{L}_{\rm res}(\theta; \alpha)
= \frac{1}{B}\sum_{b=1}^B \big(E_L(\tilde R_b)-E_{\rm eff}(\alpha)\big)^2,
\qquad
E_{\rm eff}(\alpha)=\alpha\,E_{\rm DMC}+(1-\alpha)\,\mu.
\end{equation}
\begin{itemize}
  \item \textbf{Self–target warmup ($\alpha=0$).} Minimising $\mathbb{E}[(E_L-\mu)^2]$ is variance minimisation \emph{under the sampling measure in use}; $\mu$ is treated as a detached constant, so no gradient flows through the batch mean. It teaches the network to produce smooth, low–variance amplitudes inside nodal pockets and reduces sensitivity to early sampler noise.
  \item \textbf{Anneal to reference ($\alpha\uparrow$).} As capacity grows, we blend in the external target, gently steering the mean energy while retaining the variance pressure. The cosine ramp avoids abrupt objective switches.
  \item \textbf{Reference-target residual ($\alpha=1$).} Stage I may end with $E_{\rm eff}=E_{\rm DMC}$. Despite its name in the code, this is not variational energy minimisation: it is a squared residual against a fixed scalar, and it can in principle drive $\langle E_L\rangle$ \emph{up} toward a target the ansatz could beat. The variational optimisation is stage II.
\end{itemize}
In practice we compute $\nabla_\theta \mathcal{L}_{\rm res}$ by autodiff, use micro--batches, apply optional gradient clipping, and log per–component usage and shell radii in physical units.

\paragraph{Local–energy derivatives in practice.}
We support three Laplacian backends:
\begin{enumerate}
  \item \textbf{Exact:} chunked second partials (reference–accurate, heavier).
  \item \textbf{HVP–Hutchinson:} $\Delta \log\Psi = \mathbb{E}_v[v^\top H_{\log\Psi} v]$ with Rademacher $v$ (fast and stable; fallback to FD if a probe is non–finite).
  \item \textbf{FD–Hutchinson:} centred finite differences of directional gradients (robust in corner cases; slightly noisier).
\end{enumerate}
We average a small number of probes (2–16 in ablations), and trim quantiles before forming losses.

\section{Stochastic Reconfiguration (stage II)}
\label{sec:sr-stage}

Stochastic reconfiguration~\cite{Sorella1998-SR,BeccaSorella2017-QMCBook} preconditions the
energy gradient with the (inverse) quantum geometric tensor, and is the imaginary-time /
natural-gradient step for a variational state~\cite{Amari1998-NaturalGradient,StokesEtAl2020-QuantumNaturalGradient}.
It uses the covariance form of the energy gradient
\begin{equation}
\partial_{\theta_k}E
= 2\,\mathrm{Cov}_\pi\!\big(E_L, O_k\big)
= 2\Big(\langle E_L O_k\rangle - \langle E_L\rangle\,\langle O_k\rangle\Big),
\end{equation}
and approximates a natural–gradient step by solving $(S+\lambda I)\Delta\theta=-g$ with the
quantum geometric tensor $S$ of~\S\ref{sec:theory-tangent-kernel}.
Implementation details:
\begin{itemize}
  \item \textbf{Centering/whitening.} We centre $O$ per batch and whiten by the diagonal of $S$ before CG to improve conditioning.
  \item \textbf{Trust region.} The final step is scaled to satisfy a maximum parameter norm.
  \item \textbf{Damping.} A small $\lambda$ stabilises low--variance directions.
\end{itemize}
Because our backflow only alters the Slater coordinates (and $W_\theta$ stays on safe features of $\tilde R$), the SR statistics remain well–behaved even when nodes shift.

\paragraph{Optimiser variants compared.}
Three first-order updates appear in the results, and their comparison is itself a finding
(\S\ref{sec:kernel-cond}):
\begin{itemize}
  \item \textbf{Adam}~\cite{KingmaBa2015-Adam}: a diagonal adaptive step, used as the
    robust default, with separate learning rates for the Jastrow and backflow blocks.
  \item \textbf{Plain SR}: the natural-gradient step above with a fixed diagonal damping
    $\lambda$.
  \item \textbf{CG-SR}: the same step solved matrix-free by preconditioned conjugate
    gradients, with Tikhonov damping and a trust-region cap. It is most useful where $S$ is
    anisotropic but still statistically resolved. Whether that condition holds, and what
    follows when it does not, is \S\ref{sec:kernel-cond}.
\end{itemize}
Solving the SR step without forming $S$ is what makes it affordable at
$\mathcal{O}(10^5)$ parameters. The same obstacle is addressed from the other side by the
minSR and kernel-formulation reductions, which exploit the fact that the number of samples,
rather than the number of parameters, sets the true rank of the
problem~\cite{ChenHeyl2024-minSR,Rende2024-KernelSR}. Whether the extra cost buys anything
over a diagonal preconditioner is measured in \S\ref{sec:kernel-q2}.

\section{Sampling paradigms and the weak-form collocation objective}
\label{sec:method-collocation}
The alternative pipeline described in this section draws its batches from an analytic
proposal rather than a Markov chain---the SR route of \S\ref{sec:sr-stage} samples
$|\Psi_\theta|^2$ by Metropolis---so the choice of \emph{sampling measure} must be stated
explicitly: the results treat it as a variable in its own right. Two paradigms appear in this
thesis. In \emph{variational Monte Carlo} (VMC) the batch is drawn from
$|\Psi_\theta|^2$ itself, usually by a Metropolis chain, so the estimators of the QGT
and of the gradient share the very metric they are meant to precondition. In
\emph{collocation} the batch is drawn from a fixed analytic proposal $q(\mathbf R)$ and
reweighted to $|\Psi_\theta|^2$ by importance weights
$w_i=|\Psi_\theta(\mathbf R_i)|^2/q(\mathbf R_i)$, so that the gradient training requires no
Markov chain. A Metropolis probe is used only outside the gradient loop---for periodic
checkpoint selection and for the final energy certification.

For the collocation objective we use a \emph{weak form}: the local energy $E_L$ enters
only as a reward in a score-function (REINFORCE) gradient~\cite{Williams1992-REINFORCE},
\begin{equation}
  \nabla_\theta \langle E\rangle
  = 2\,\mathbb{E}_{\mathbf R\sim|\Psi_\theta|^2}\!\big[(E_L(\mathbf R)-b)\,
    \nabla_\theta \log|\Psi_\theta(\mathbf R)|\big],
  \label{eq:reinforce-method}
\end{equation}
with a variance-reducing baseline $b$ (a running mean of $E_L$), and the Laplacian inside
$E_L$ is never differentiated through. This is deliberate. Differentiating the Laplacian
through a backflow-shifted determinant is the source of the nodal instability analysed in
Appendix~\ref{app:catch22}; the weak form keeps the full local-energy signal as a reward
while sidestepping that pathology. Two robustness controls stabilise the reward: local
energies are clipped to a robust median/MAD location--scale band ($\sim\!5\sigma$), and the
most extreme rewards are quantile-trimmed before the gradient is formed.

\subsection{Analytic proposal and importance resampling}
\label{subsec:colloc-proposal}
The collocation points are drawn from an analytic mixture of isotropic Gaussians centred at
the trap minimum and reweighted to $|\Psi_\theta|^2$~\cite{KalosLevesqueVerlet1974-ImportanceSampling},
\begin{equation}
  q(\mathbf R)=\frac{1}{K}\sum_{k=1}^{K}\mathcal N\!\big(\mathbf R;\,0,\,\sigma_{f,k}^2/\omega\;I_{Nd}\big),
  \quad \sigma_f\in\{0.8,1.3,2.0\},
  \qquad
  w_i=\frac{|\Psi_\theta(\mathbf R_i)|^2}{q(\mathbf R_i)},
  \label{eq:colloc-proposal}
\end{equation}
where the $1/\omega$ scaling ties the proposal width to the oscillator length
$\ell=1/\sqrt\omega$ and the three components cover core, bulk, and tails. A batch of
$n_{\rm coll}$ points is obtained by multinomial resampling on $p_i=w_i/\sum_j w_j$ from
$m\,n_{\rm coll}$ candidates ($m\!\approx\!8$--$16$). Reweighting quality is tracked by the
effective sample size $\mathrm{ESS}=(\sum_i w_i)^2/\sum_i w_i^2$ and, as a heavier-tailed
diagnostic, the Pareto shape parameter $\hat k$ of the
weights~\cite{Vehtari2024}. The second is not a second opinion on the first. Importance
weights are heavy-tailed, and $\hat k$ estimates the tail index of the generalised Pareto
distribution fitted to the largest of them. The thresholds are statements about the
moments of the \emph{fitted} tail. For $\hat k<1/2$ the weights behave as though the
variance were finite and the reweighted estimator obeys a central limit theorem. Between
$1/2$ and $1$ the variance behaves as infinite though the mean does not, and convergence is
slow enough that $\hat k\approx0.7$ is the practical limit of trustworthiness. Above $1$ the
fitted tail has no finite mean.

That last statement needs care, because in this setting the underlying ratio does have a
finite mean by construction: $w=|\Psi_\theta|^2/q$ with $q$ a proper density gives
$\mathbb{E}_q[w]=\int|\Psi_\theta|^2<\infty$ whenever $\Psi_\theta$ is square-integrable, so
asymptotically $\hat k<1$. A measured $\hat k>1$ is therefore a \emph{finite-sample}
diagnostic: it says the sample is deep in a pre-asymptotic regime where the realised
estimator behaves as if no finite mean existed, not that the expectation fails to exist
\cite{Vehtari2024}. We use it in that sense throughout. An ESS can look survivable while $\hat k$ says the
estimator has no variance to speak of, which is why we track both. Both degrade when $q$
ceases to overlap $|\Psi_\theta|^2$, and that degradation is the failure mode which bounds
the paradigm at large $N$ and weak $\omega$ (Chapter~\ref{ch:kernel},
\S\ref{sec:colloc-reliability}). Three controls manage the tails: weight tempering ($w\!\to\!w^{\alpha}$, $\alpha<1$),
log-weight clipping at an upper quantile, and ESS-adaptive oversampling that raises $m$ when
the achieved ESS falls below a floor. The first two are deliberate bias--variance trades:
they change the distribution the reward is averaged over, so the reweighted energy readout is
no longer an estimate of $\langle E_L\rangle_{|\Psi_\theta|^2}$. That is why the final
certification is always a separate evaluation under $|\Psi_\theta|^2$
(\S\ref{sec:kernel-cond}). An instability-rollback reverts to the
last good checkpoint and decays the learning rate if the energy jumps beyond a threshold.

\subsection{Physics-informed proposal for the Wigner regime}
\label{subsec:colloc-wigner}
The origin-centred proposal~\eqref{eq:colloc-proposal} places its mass where a Wigner
molecule has almost none: at strong correlation the density collapses onto a shell at the
classical ring radius $r_0$, and $|\Psi_\theta|^2$ has essentially no support near the
origin, so the ESS collapses. We therefore replace $q$ by a proposal built from the Wigner
shell itself---a radial Gaussian centred at $r_0$ (set by balancing confinement against
Coulomb repulsion) times an angular distribution over the ring, so that proposed
configurations resemble electrons spread around the ring instead of clustered at the centre.
Where the ring itself develops angular structure the uniform angular factor is replaced by a
Dirichlet model of the angular gaps. Difficult confinements are reached by \emph{cascade
warm-starting}: parameters converged at one $\omega$ initialise the next, following
$\omega:1.0\!\to\!0.5\!\to\!0.1\!\to\!0.01\!\to\!0.001$. What this construction achieves,
and where it stops, is \S\ref{sec:kernel-q3-frontier}.

\subsection{Ansatz families used in the collocation programme}
\label{subsec:colloc-families}
Three ansatz families are trained under this objective in Chapter~\ref{ch:kernel}, all on
the same Slater--Jastrow core $\Psi(\mathbf R)=\det[\phi_i(\mathbf r_j+\Delta\mathbf r_j)]\,e^{J(\mathbf R)}$
and differing only in what dresses it. They are deliberately smaller than the production
networks of \S\ref{sec:hyperparameters}, because this programme is limited by gradient noise and memory, not by capacity.

\emph{Jastrow-only.} The correlator alone at ${\sim}25$k parameters, in its
message-passing form (\texttt{CTNN\-Jastrow\-VCycle}), acting on node features
(single-particle orbital values) and edge features (soft-core pair distances). It is used
both standalone and as the initialiser for the backflow runs.

\emph{Backflow $+$ Jastrow.} Adds the copresheaf displacement network
(\texttt{CTNNBackflowNet}), ${\sim}49$k parameters in total at $h_{\rm bf}{=}128$. Its edge
tensor scales as $\mathcal O(N^2 h_{\rm bf})$, which is the memory bottleneck at large $N$:
at $N{=}20$ the hidden width must be cut to $h_{\rm bf}\in\{48,64\}$.

\emph{Neural Pfaffian.} Replaces the determinant with $\mathrm{Pf}[A_{ij}]$, a more general
antisymmetric form admitting richer nodal topology~\cite{GaoGuennemann2024-NeuralPfaffian}.

\section{Units, scaling, and batch construction}
\label{sec:units-scaling}

All features are computed in trap units $\tilde{\mathbf x}=\sqrt{\omega}\,\mathbf x$ so that length scales and gates (e.g.\ the short–range gate in the NQS and backflow) behave consistently across $\omega$. The sampler internally widens/narrows Gaussians and shell radii using $a_{\rm ho}=1/\sqrt{\omega}$ in a fixed, explicit way, so that the five components cover comparable \emph{dimensionless} regimes for $\omega\in[0.01,1]$. Proposal candidates are drawn independently, so the batch carries no serial Markov-chain autocorrelation; the retained batch is not i.i.d.\ in the strict sense, since multinomial resampling can repeat candidates and the hard-example buffer reinjects earlier ones. What the design buys is the absence of a mixing time, not statistical independence.

\section{The two workflows}
\label{sec:opt-summary}

Both routes start from the same ansatz and end at an energy evaluated under
$|\Psi_\theta|^2$. They differ in the measure that defines the loss in between.
\begin{center}\small
\begin{tabular}{@{}ll@{}}
\textbf{SR--VMC (\S\ref{sec:residual-stage}--\S\ref{sec:sr-stage})} &
\textbf{Collocation (\S\ref{sec:method-collocation})}\\[2pt]
stratified mixture $\to$ residual pretraining &
pretrained checkpoint\\
$\downarrow$ & $\downarrow$\\
$|\Psi_\theta|^2$ sampling (Metropolis) &
analytic proposal $q$, importance reweighting\\
$\downarrow$ & $\downarrow$\\
SR step $(S+\lambda I)\Delta\theta=-g$ &
weak-form REINFORCE gradient (Adam)\\
$\downarrow$ & $\downarrow$\\
energy under $|\Psi_\theta|^2$ &
Metropolis probe for checkpoint choice,\\
 & then energy under $|\Psi_\theta|^2$\\
\end{tabular}
\end{center}
The right-hand route places no Markov chain in any gradient; the probe that selects its
checkpoints sits outside the update loop entirely. The left-hand route, which carries the
ground-state energies of \S\ref{sec:energy-overview}, keeps the early dynamics
variance-dominated and sampler-aware and then transitions to a geometry-aware
natural-gradient refinement---a division that matters most at weak traps, where the shell and
dimer coverage of the stratified mixture is what makes the early stage converge at all.

\section{Statistical uncertainties}
\label{sec:uncertainties}
Every energy and observable we report carries a statistical error, and two distinct sources
contribute. Within a single evaluation, successive MCMC samples are autocorrelated, so a
naive standard error over $B$ samples understates the sampling uncertainty; this is what
blocking addresses. Across training runs, the optimisation is itself stochastic and lands on
different states from different seeds; blocking says nothing about that, and it is assessed
separately by repeating runs at independent seeds wherever the claim depends on it
(\S\ref{sec:colloc-reliability}, \S\ref{sec:kernel-q1a}). We therefore
estimate the error on a scalar mean $\bar A=\frac1B\sum_i A(\mathbf R_i)$ by
\emph{blocking}: successive samples are averaged into blocks of growing size and the
plateau of the block standard error as the block size exceeds the autocorrelation time is
taken as the uncertainty~\cite{FlyvbjergPetersen1989-Blocking,Jonsson2018-BlockingSE}. For
the final ground-state energy this is the dominant reported error; a parenthetical digit,
e.g.\ $155.8738(7)$, denotes one standard error on the last quoted figure.

For quantities that are \emph{nonlinear} functions of sample statistics (participation ratios, effective ranks, PCA cosines, and the coupling diagnostics of Chapter~\ref{ch:results}) a blocked standard error is not available in closed form, and we
use a \emph{block bootstrap} instead: blocks (of the plateau size above) are resampled with
replacement, the statistic is recomputed on each resample, and its spread across resamples
is the reported error~(\S\ref{subsec:repr-uncertainty}). Where a number is quoted as a
seed average, we additionally report the spread across independent random seeds, which captures optimisation variability and not only sampling variability; the two are distinguished in the
captions. Reference energies drawn from the literature carry the uncertainties quoted by
their sources.

\chapter{Analysis}
\label{ch:analysis}

This chapter describes \emph{what} we measure on fresh $|\Psi_\theta|^2$ samples, \emph{how} we compute the diagnostics, and \emph{why} each quantity is informative. Results and numbers are presented in Chapter~\ref{ch:results}.
\section{Structural diagnostics and Wigner markers}
\label{sec:analysis-struct}

Unless stated otherwise, curves are accumulated in \emph{trap units} 
$\mathbf{y}=\sqrt{\omega}\,\mathbf{x}$ to compare shapes across $\omega$, 
while all headline distances are reported in \emph{Bohr} (physical units $\mathbf{x}$).
Production chains are thinned; we split samples into halves for stability checks 
(acceptance window, JSD between halves, and reproducibility under restarts).
Using trap units to compare shapes across confinements is standard in parabolic dots; 
see the overview in~\cite{manninen2007metalclustersquantumdots}.

\subsection{Pair distribution and radial statistics}
\label{subsec:gr-pr}

We estimate an \emph{effective} pair distribution---a shell-area-corrected pair-distance
density---in trap units. For a finite trapped dot there is no homogeneous-system
normalisation, because the one-body density varies strongly across the cloud, so $A_{\rm eff}$
below is a finite-dot convention rather than a thermodynamic volume and $g(r)$ does not tend
to unity at large $r$. Peak positions and relative profiles, which is what the analysis uses,
are insensitive to that choice; absolute amplitudes are not.
\begin{equation}
g(r_\alpha)\;=\;\frac{A_{\mathrm{eff}}}{N(N-1)}\;
\frac{\big\langle\sum_{i<j}\mathbf{1}\{r_{ij}\!\in[r_\alpha^-,r_\alpha^+)\}\big\rangle}
{2\pi r_\alpha\,\Delta r_\alpha},
\qquad r_{ij}=\|\mathbf{y}_i-\mathbf{y}_j\|_2,
\end{equation}
with $A_{\mathrm{eff}}=\pi R_{0.95}^2$ and $R_{0.95}$ the 95\% quantile of single‐particle radii in trap units.
This $g(r)$-based structural analysis follows established practice in mesoscopic 2D electron systems
and quantum dots~\cite{Filinov_2001,Egger_1999,manninen2007metalclustersquantumdots}.
The pair-distance probability is therefore denoted
\(
P_{\rm pair}(r_{ij})\propto r_{ij}^{d-1}g(r_{ij})
\).
It is used for pair-structure diagnostics, distinct from the one-body radial
probability
\(
P_1(r)\propto r\rho_1(r)
\),
from which we report the headline mode $r_{\mathrm{mode}}$, mean $\bar r$,
standard deviation $\sigma_r$, and FWHM (all in Bohr). A compact one-body
width-to-scale marker is
\begin{equation}
\gamma \;=\; \sigma_r / r_{\mathrm{mode}},
\label{eq:gamma}
\end{equation}
which tracks the width of $P_1$ relative to its modal radius. It is interpreted
alongside shell separation and pair/angular diagnostics, rather than as a
universal monotone Wigner marker~\cite{Egger_1999,Filinov_2001}.

\subsection{Automatic shell detection and topologies}
\label{subsec:shells}

For each frame we sort the single‐particle radii $s_1\le\cdots\le s_N$ (trap units), compute gaps
$g_k=s_{k+1}-s_k$, and flag the frame as \emph{two‐shell} if the largest gap exceeds a robust threshold:
\begin{equation}
\max_k g_k \;\ge\; \tau\,\mathrm{median}(g_k),
\qquad \tau\in[2.5,3.5]\ \text{(default 3.0)}.
\end{equation}
The cut radius is the midpoint of the maximal gap; the inner multiplicity $n_{\mathrm{in}}$ defines the topology 
$(n_{\mathrm{in}}, N-n_{\mathrm{in}})$.

\paragraph{More than two shells.} At $N{=}20$ and weak confinement a single cut is not
enough, and the detector is applied recursively. Having placed one cut, we re-apply the same
robust-gap test \emph{within} each resulting group, accepting a further cut only if that
group contains at least three particles and its own maximal gap again exceeds
$\tau\,\mathrm{median}(g_k)$ computed over the whole frame---so the threshold is global and
a shell is never split by a gap that is small on the scale of the configuration. Recursion
stops when no group qualifies, which for the systems here never yields more than three
shells. A frame is labelled by the resulting occupancy vector
$(n_1,\dots,n_S)$ ordered outward, so that the two-shell case above is the $S{=}2$ member of
the same family; the $N{=}20$ deep-Wigner topology $(1,7,12)$ is an $S{=}3$ assignment with
cuts after the first and eighth particle. Topology fractions and their robustness under
$\tau$ are reported exactly as in the two-shell case.
Shelling and discrete ring occupancies are well known in parabolic Coulomb clusters
and quantum dots~\cite{schweigert1994spectralpropertieschargedparticles,Kong_2002,manninen2007metalclustersquantumdots};
our detector provides a reproducible, data-driven assignment in that spirit.
We report:
(i) the fraction of two‐shell frames,
(ii) the histogram over $n_{\mathrm{in}}$ among two‐shell frames,
and (iii) for selected topologies (e.g.\ $(1,5)$ or $(3,9)$) the fraction among \emph{all} frames.

\paragraph{Shell‐resolved reconstruction of $g(r)$.}
On two‐shell frames we decompose pairs into II, IO, and OO and accumulate histograms
$g_{\mathrm{II}},g_{\mathrm{IO}},g_{\mathrm{OO}}$. With empirical weights
$w_{\mathrm{II}},w_{\mathrm{IO}},w_{\mathrm{OO}}$ (combinatorial counts) we reconstruct
\(\tilde g(r)=w_{\mathrm{II}}g_{\mathrm{II}}+w_{\mathrm{IO}}g_{\mathrm{IO}}+w_{\mathrm{OO}}g_{\mathrm{OO}}\)
and report $\cos\angle(g,\tilde g)$.
Values $\approx 1$ indicate that the observed $g(r)$ shape is explained by shell geometry, not sampling artefacts,
in line with shell-based descriptions in~\cite{schweigert1994spectralpropertieschargedparticles,Kong_2002,manninen2007metalclustersquantumdots}.

\subsection{Orientational order and “ring stiffness”}
\label{subsec:orient}

On each detected ring (inner or outer) we take the particle angles $\{\phi_k\}_{k=1}^n$ in
the lab frame. No rotation is removed at this stage: the modulus $|\Phi_m|$ below is already
invariant under a global rotation, and the registration of Eq.~\eqref{eq:Ctilde} is applied
only where the \emph{phase} is needed, which is the cross-shell comparison.
Bond‐orientational order is
\begin{equation}
\Phi_m \;=\; \frac{1}{n}\sum_{k=1}^n e^{\,\mathrm{i} m \phi_k},
\qquad m=n,
\end{equation}
so that $|\Phi_m|\in[0,1]$ measures $n$‐fold orientational order (e.g.\ $m{=}5$ for $(1,5)$, $m{=}9$ for the outer ring of $(3,9)$),
analogous to standard bond–orientational order parameters used in 2D Coulomb/Yukawa crystals~\cite{Mazars_2008}.
As a complementary, dimensionless \emph{angular Lindemann} number we use
\begin{equation}
\lambda_\phi \;=\; \frac{\mathrm{std}(\Delta\phi)}{\mathbb{E}[\Delta\phi]},
\qquad 
\Delta\phi\;=\;\text{nearest‐neighbour angular spacing on the ring}.
\end{equation}
Smaller $\lambda_\phi$ indicates a stiffer (more Wigner‐like) ring. This is the quantity
abbreviated ``Lind'' in Chapter~\ref{ch:results}; where a shell index is attached, as in
$\lambda_{\rm in}$ and $\lambda_{\rm out}$, it is the same estimator evaluated on that
shell alone.

\paragraph{Registered and lab-frame order.}
Two further scalars are used for the multi-shell cases, and both are needed because
$|\Phi_{n_s}|$ is defined shell by shell while the question at $N{=}20$ is whether the
\emph{molecule} is ordered. For a frame with shells $s=1,\dots,S$ and occupancies $n_s$ we
define a configuration-level order parameter in two frames. In the \emph{lab} frame,
\begin{equation}
  C_{\rm global} \;=\; \Big|\;\Big\langle \tfrac{1}{S'}\textstyle\sum_{s:\,n_s\ge3}
  \Phi_{n_s}\Big\rangle_{\rm frames}\Big| ,
  \label{eq:Cglobal}
\end{equation}
the modulus of the \emph{frame-averaged} complex order parameter, averaged over the
$S'$ shells carrying at least three particles. Because a freely rotating molecule visits all
orientations, the complex phases cancel and $C_{\rm global}\to0$ however ordered the
molecule is internally: it measures \emph{pinning} to the lab frame, not order. In the \emph{co-rotating}
frame each configuration is first registered by an estimated physical rotation. Because a
rigid rotation by $\alpha$ sends $\Phi_{n_s}\mapsto e^{in_s\alpha}\Phi_{n_s}$, the phase of a
shell's order parameter advances at $n_s$ times the rotation angle, and the angle itself must
be recovered from the outermost shell as
$\hat\alpha=\arg\Phi_{n_S}/n_S$ (modulo $2\pi/n_S$) before it can be applied to the others:
\begin{equation}
  \tilde C_{\rm global} \;=\; \Big|\;\Big\langle \tfrac{1}{S'}\textstyle\sum_{s:\,n_s\ge3}
  \Phi_{n_s}\,e^{-i n_s\hat\alpha}\Big\rangle_{\rm frames}\Big| ,
  \qquad \hat\alpha=\frac{\arg\Phi_{n_S}}{n_S}.
  \label{eq:Ctilde}
\end{equation}
Dividing the outer phase by $n_S$ is essential: removing $\arg\Phi_{n_S}$ from every shell
alike would only correspond to a common spatial rotation when all occupancies coincide, and
would otherwise mix shells that rotate together into an apparent misalignment.
This retains the \emph{relative} angular arrangement of the shells while discarding the
global orientation, so a rigid molecule tumbling freely gives $\tilde C_{\rm global}$ near
its single-frame value and $C_{\rm global}$ near zero. The pair
$(\tilde C_{\rm global}, C_{\rm global})$ is therefore the direct test of the rotating
Wigner-molecule picture, and it is reported that way in \S\ref{sec:wigner-angular}.
Finally, $L_{n,\rm reg}$ denotes the angular Lindemann number $\lambda_\phi$ evaluated on
the shell of occupancy $n$ after the same registration.

Two of these quantities are insensitive to the registration convention and one is not, and
the distinction bounds what may be concluded from each. The shellwise modulus
$|\Phi_{n_s}|$ and the nearest-neighbour spacing statistics behind $\lambda_\phi$ are
invariant under any global rotation, so they carry no dependence on how $\hat\alpha$ is
estimated. The \emph{cross-shell} quantity $\tilde C_{\rm global}$ does depend on it, since
its whole content is the relative alignment between shells of different occupancy. The
$N{=}20$ values reported in \S\ref{sec:wigner-angular} were produced by an earlier analysis
pass whose code has not been retained, so we cannot confirm which convention it used; the
per-shell numbers there are unaffected, but the cross-shell $\tilde C_{\rm global}$ should be
recomputed under Eq.~\eqref{eq:Ctilde} before it is relied on quantitatively.

\paragraph{The random-angle null.}
None of these is interpretable without a baseline. For $n$ particles at independent uniform
angles, $n\Phi_n$ is a two-dimensional random walk of $n$ unit steps, so
$\langle|\Phi_n|^2\rangle=1/n$ exactly, and the mean modulus follows from Kluyver's integral,
\begin{equation}
  \big\langle|\Phi_n|\big\rangle_{\rm random}
  =\frac1n\int_0^\infty\frac{1-J_0(t)^n}{t^2}\,dt ,
  \label{eq:phi-null-exact}
\end{equation}
which we evaluate numerically. The familiar Rayleigh form $\sqrt\pi/(2\sqrt n)\approx0.886/\sqrt n$
is its large-$n$ limit and undershoots it at the ring sizes that occur here, by $2.6\%$ at
$n=3$, $1.3\%$ at $n=5$ and $0.4\%$ at $n=19$---the same size as the effects being tested at
strong confinement, which is why the exact value is used throughout. Because shells are
assigned by radial rank alone and no angular cut is applied, it is also exactly the null of
our pipeline: redrawing every angle uniformly while keeping each frame's radii reproduces it
within its bootstrap error in all fifteen $(N,\omega)$ cells
(Appendix~\ref{app:structural-tables:shells}). A five-particle ring with no angular order
whatever returns $|\Phi_5|\approx0.40$. We therefore report $|\Phi_n|$ together with its ratio to this null
wherever the ordering claim depends on it (Appendix~\ref{app:structural-tables:shells}).
This is a Lindemann–type orientational metric tailored to finite rings, motivated by
the use of Lindemann criteria to diagnose intershell rotational (angular) melting in mesoscopic
Coulomb clusters and dots~\cite{schweigert1994spectralpropertieschargedparticles,Filinov_2001}.

\subsection{Scaling check for $N=2$}
\label{subsec:n2-scaling}

To connect with the classical strong‐coupling limit we fit 
\begin{equation}
r_{\mathrm{mode}}(\omega) \;=\; C\,\omega^{\alpha}
\end{equation}
by linear regression in $\big(\log\omega,\log r_{\mathrm{mode}}\big)$ and compare the slope with the classical reference 
$\alpha_{\mathrm{cl}}=-2/3$ (two charges in a harmonic trap at strong coupling; see, e.g., the reviews in~\cite{manninen2007metalclustersquantumdots}).
Approach toward $\alpha_{\mathrm{cl}}$ as $\omega\downarrow$ is one of our Wigner crossover indicators,
consistent with the crossover phenomenology reported in~\cite{Egger_1999,Filinov_2001}.

\subsection{Reporting and stability checks}
\label{subsec:analysis-sanity}

For all scalars and histograms we report standard errors calculated over thinned blocks or frames.
We verify: 
(i) acceptance in the target window, 
(ii) Jensen–Shannon divergence and $\ell_2$ distance between halves of production (stationarity),
(iii) robustness of topology fractions and $|\Phi_m|$ under $\tau\in[2.5,3.5]$ and mild pre‐smoothing of radii, and
(iv) reproducibility of $g(r)$ and $P(r)$ under independent sampler restarts.
All diagnostics operate directly on saved bundles 
$\{\mathbf{X}_{\mathrm{Bohr}},\,r,\,g(r)\}$.

\section{Representation analysis}
\label{sec:repr-methods}

We probe what the networks learn by evaluating fixed, trained models on fresh
$|\Psi_\theta|^2$ samples. Unless stated otherwise, features are computed in \emph{trap units}
$\mathbf{y}=\sqrt{\omega}\,\mathbf{x}$ to match the correlator’s internal scaling.
All estimates use thinned MCMC frames, split into two equal blocks $A$ and $B$;
PCA fits on $A$ and reports on $B$ (and vice versa) to avoid circularity. 
We use \texttt{float64} throughout.

\paragraph{Feature taps.}
Let $z_\theta\in\mathbb{R}^D$ denote the concatenated correlator representation
\[
z_\theta=\big[\overline{\phi}\;\big|\;\overline{\psi}\;\big|\;\mathbf{g}\big],
\]
where $\overline{\phi}$ are per-particle summaries, $\overline{\psi}$ are pairwise summaries,
and $\mathbf{g}$ are global/extras. The readout $\rho_\theta$ is a two-layer MLP, not a linear
map, so
\begin{equation}
f_{\rm base}(z)=\rho_\theta(z),
\qquad
f=f_{\rm base}+f_{\rm cusp},
\end{equation}
and the ablations below re-evaluate $\rho_\theta$ on modified inputs instead of manipulating a weight vector.
Backflow returns a displacement field $\Delta x\in\mathbb{R}^{N\times d}$ that enters only the
Slater determinant; the correlator and cusp are evaluated on the unshifted $\tilde R$
(\S\ref{subsec:backflow-method}). The coupling diagnostics below additionally probe the
correlator on the shifted coordinates as a counterfactual sensitivity.

\subsection{Representation geometry}
\paragraph{Preprocessing and PCA.}
\label{subsec:repr-prep}
Given a sample matrix $Z\in\mathbb{R}^{K\times D}$ (rows are frames), we \emph{mean-centre}
the columns, $Z_c=Z-\mu$, and take the SVD of $Z_c$ to obtain orthonormal loadings
$U=[u_1,\dots,u_D]$; scores are $S=Z_cU$. Centring but not rescaling is deliberate: the
readout $\rho_\theta$ was trained on the raw features, so every ablation below is performed
in raw coordinates and the mean is restored before $\rho_\theta$ is applied. Where a
per-column standardisation is used---for the branch-power decomposition of
\S\ref{subsec:repr-pc1-power}, which compares loadings across branches carrying different
natural scales---it is stated explicitly. PCA is fitted on block $A$ and reported on the
held-out block $B$ unless noted.

\paragraph{Entropy effective rank.}
\label{subsec:repr-erank}
Let $\{\lambda_i\}_{i=1}^D$ be the eigenvalues of $\mathrm{Cov}(Z_c)$
(on the held-out block). With $p_i=\lambda_i/\sum_j\lambda_j$, the
\emph{entropy effective rank} is
\[
r_{\rm eff}(Z)\;=\;\exp\!\Big(-\sum_{i=1}^D p_i\log p_i\Big),
\]
which equals $1$ for perfectly one-dimensional representations and increases
with spectral spread. Two spectral counts appear in this thesis and they are not the same
number. $r_{\rm eff}=\exp(H_1)$ above is the exponential of the Shannon entropy of the
spectrum; the \emph{participation ratio}
$d_{\rm eff}=(\sum_a\lambda_a)^2/\sum_a\lambda_a^2$ of \S\ref{sec:theory-tangent-kernel} is
the exponential of the R\'enyi-2 entropy. Since R\'enyi entropies decrease in their order,
$d_{\rm eff}\le r_{\rm eff}$ always, with equality only for a flat spectrum. Appendix~\ref{app:representation}
reports $r_{\rm eff}$ for feature covariances and Chapter~\ref{ch:kernel}
reports $d_{\rm eff}$ for the tangent metric and for the displacement field, so the same
object can carry two slightly different counts in the two chapters; each table names which
is used.

\subsection{Probes and ablations}
\paragraph{Principal-component ablations of the head.}
\label{subsec:repr-pc-ablate}
To test how many latent axes the readout truly uses, we project the centred features onto
the top-$k$ PCs and restore the mean before re-evaluating the trained readout \emph{without
refitting},
\begin{equation}
\Pi_k=\sum_{i=1}^k u_i u_i^\top,
\qquad
z^{(k)}=\mu+\Pi_k\,(z-\mu),
\qquad
\hat f^{(k)}=\rho_\theta\big(z^{(k)}\big).
\label{eq:pc-ablation}
\end{equation}
Restoring $\mu$ matters: $\rho_\theta$ is a nonlinear map trained on uncentred inputs, so
feeding it a centred or standardised vector would evaluate it off its training domain and
the resulting error would measure that displacement rather than the loss of PC content. As a
null, the same projection is repeated onto $k$ \emph{random} orthonormal directions and the
two error curves are compared; a top-$k$ curve that is not below the random-$k$ curve would
mean the leading PCs carry no privileged information for the readout.
We report a \emph{normalised} error on block $B$,
\[
\mathrm{rel\text{-}MAE}(k)\;=\;
\frac{\mathbb{E}_B\,|\,\hat f^{(k)}-f_{\rm base}\,|}
     {\mathrm{std}_B(f_{\rm base})},
\]
together with the learning curve in $k$.

\paragraph{PC1 block-power decomposition.}
\label{subsec:repr-pc1-power}
To attribute the leading axis to branches, we partition $u_1$ conformably with
$z=[\overline{\phi}\,|\,\overline{\psi}\,|\,\mathbf{g}]$ and report
\[
\text{power}(\overline{\phi})=\frac{\|u_{1,\overline{\phi}}\|_2^2}{\|u_1\|_2^2},
\quad
\text{power}(\overline{\psi})=\frac{\|u_{1,\overline{\psi}}\|_2^2}{\|u_1\|_2^2},
\quad
\text{power}(\mathbf{g})=\frac{\|u_{1,\mathbf{g}}\|_2^2}{\|u_1\|_2^2}.
\]
These are computed on standardised features so that branches carrying different natural
scales are comparable. One confound remains and should be kept in mind when reading them: a
branch holding $d_b$ of the $D$ coordinates would already carry about $d_b/D$ of the squared
loading mass under an isotropic random direction, so a large share partly reflects a wide
branch. The shares below are therefore read as \emph{relative} shifts across $\omega$ at
fixed architecture---where the branch widths are constant and the confound cancels---and not
as absolute statements that one branch dominates.

\paragraph{Linear probes to coarse physics.}
\label{subsec:repr-probes}
We regress interpretable summaries per configuration on the first $k$ PC scores $S^{(k)}$ using
ridge regression (tiny penalty $\alpha=10^{-6}$):
\[
y \approx S^{(k)} \beta,\qquad 
y\in\big\{\bar r,\ \mathrm{Var}(r),\ \Pr(r<r_0),\ \text{shell\ contrast}\big\}.
\]
We report $R^2$ on the held-out block; $r_0$ is a fixed near-contact cutoff in trap units.

\paragraph{Near-field gradient concentration.}
\label{subsec:repr-nearfield}
For each frame we compute the minimum pair distance $r_{\min}$ (trap units) and the gradient
of the residual head with respect to particle coordinates, $\nabla f_{\rm base}(\mathbf{y})$.
The quantity reported is an \emph{enrichment factor}: the mean gradient power among the
closest-contact frames divided by the mean over all frames,
\begin{equation}
\mathcal{E}(q)\;=\;
\frac{\mathbb{E}\big[\,\|\nabla f_{\rm base}\|_2^2 \;\big|\; r_{\min}\le Q_q\,\big]}
     {\mathbb{E}\big[\|\nabla f_{\rm base}\|_2^2\big]},
\label{eq:nf-enrichment}
\end{equation}
with $Q_q$ the $q$-quantile of $r_{\min}$. This is a ratio of a \emph{conditional} mean to an
unconditional one and is therefore unbounded above, unlike an indicator-weighted mass
fraction, which could not exceed unity. $\mathcal{E}\approx1$ means gradient power is spread
in proportion to frame count, so the analytic cusp is handling most short-range structure;
$\mathcal{E}\gg1$ means residual stiffness is concentrated near contact.

\subsection{Backflow--correlator coupling}
\paragraph{Backflow geometry and energy-locality shares.}
\label{subsec:repr-bf}

\paragraph{Field spectrum.}
We vectorize the backflow displacement per frame,
$\delta=\mathrm{vec}(\Delta x)\in\mathbb{R}^{Nd}$, standardize across frames,
and compute the entropy effective rank $r_{\rm eff}(\Delta x)$ of
Sec.~\ref{subsec:repr-erank}. Chapter~\ref{ch:kernel} reports the participation ratio of
the same spectrum instead, which is the smaller of the two counts; the ceiling
$2N-d$ set by the zero-mean projection~\eqref{eq:bf-com} is common to both.

\paragraph{Where backflow acts.}
On the same thinned frames we evaluate local energies with and without backflow,
$E_{\rm loc}^{\rm BF}$ and $E_{\rm loc}^{\rm noBF}$, using the \emph{same} coordinates.
We form the signed correction $\Delta E_{\rm loc}=E_{\rm loc}^{\rm BF}-E_{\rm loc}^{\rm noBF}$ and report
near-field shares at quantile $q$ of $r_{\min}$:
\[
\mathrm{NF}\text{-}\Delta E(q)\;=\;
\frac{\mathbb{E}\big[|\Delta E_{\rm loc}|\,\mathbf{1}\{r_{\min}\le Q_q\}\big]}
     {\mathbb{E}\big[|\Delta E_{\rm loc}|\big]}.
\]
(Using absolute values isolates magnitude; signed averages are also available.)

\paragraph{Alignment of correlator representations.}
\label{subsec:repr-alignment}
We collect two standardized feature matrices $\tilde Z^{\rm noBF}$ and $\tilde Z^{\rm BF}$
on identical frames (each standardized by its own $\mu,\sigma$), compute their PC1 loadings
$u_1^{\rm noBF}$ and $u_1^{\rm BF}$, and report the cosine similarity
\[
\cos\big(u_1^{\rm noBF},u_1^{\rm BF}\big)\;=\;
\frac{\langle u_1^{\rm noBF},\,u_1^{\rm BF}\rangle}
     {\|u_1^{\rm noBF}\|_2\,\|u_1^{\rm BF}\|_2}.
\]
Cosines near $1$ indicate that backflow acts as a structured reweighting along the same dominant axis,
rather than rotating the correlator’s manifold.

\paragraph{Centred kernel alignment.}
The PC1 cosine compares one direction. To compare the two feature \emph{geometries} as a
whole we use linear centred kernel alignment~\cite{Kornblith2019-CKA}. With
$\tilde Z^{\rm noBF},\tilde Z^{\rm BF}\in\mathbb{R}^{K\times D}$ column-centred on the same
$K$ frames,
\begin{equation}
  \mathrm{CKA}\big(\tilde Z^{\rm noBF},\tilde Z^{\rm BF}\big)
  \;=\;
  \frac{\big\|\,(\tilde Z^{\rm noBF})^{\!\top}\tilde Z^{\rm BF}\,\big\|_F^{2}}
       {\big\|(\tilde Z^{\rm noBF})^{\!\top}\tilde Z^{\rm noBF}\big\|_F\;
        \big\|(\tilde Z^{\rm BF})^{\!\top}\tilde Z^{\rm BF}\big\|_F} \;\in[0,1],
  \label{eq:cka}
\end{equation}
which is invariant to orthogonal transformations and to isotropic scaling of either
representation, and equals $1$ only when the two Gram matrices are proportional. It is the
quantity reported in Table~\ref{tab:backflow-geom}. Two cautions attach to it. CKA compares
the \emph{geometry of the feature cloud}, so a high value says the correlator's
representation is unchanged by the presence of backflow---not that the backflow occupies the
same subspace, and not that the two components are functionally independent. And CKA is
known to saturate near $1$ whenever both representations are dominated by a few shared
high-variance directions, which is exactly the regime here
(\S\ref{subsec:repr-erank}); values above $0.99$ should therefore be read as ``no detectable
reorganisation'' rather than as a quantitative distance.

\paragraph{Uncertainty and stability.}
\label{subsec:repr-uncertainty}
For scalars we report standard errors over blocks.
For PCA-based quantities and cosines we use a block bootstrap (resampling blocks with replacement)
to obtain 68\%/95\% CIs. Stability checks include:
(i) swapping train/held-out blocks for PCA,
(ii) varying the near-field quantiles $q\in\{1,5,10\}\%$,
(iii) re-running analyses on independent sampler restarts.

\paragraph{Reproducibility package.}
All diagnostics consume the saved bundles
$\{\mathbf{X}_{\rm Bohr},\,r,\,g(r)\}$ and the model snapshot,
and emit JSON/CSV tables with $r_{\rm eff}$, rel-MAE$(k)$ curves, block-power shares,
$R^2$ probes, near-field shares, $r_{\rm eff}(\Delta x)$, and noBF/BF cosines.
As in Sec.~\ref{sec:analysis-struct}, we log JSD/$\ell_2$ splits to confirm stationarity.

\section{Tangent-kernel diagnostics}
\label{sec:analysis-kernel}
The representation diagnostics above describe the ansatz's \emph{features}; a second
family, developed with the geometry of \S\ref{sec:theory-tangent-kernel}, describes its
\emph{tangent space}, and it is these that carry the architecture, optimiser, and
paradigm comparisons of Chapter~\ref{ch:kernel}. From a batch we assemble the centred
score matrix $O_{ki}=\partial_{\theta_i}\log|\Psi_\theta(\mathbf x_k)|$ and form the
quantum geometric tensor $S=O^{\!\top}O/B$ from the \emph{centred} $O$ (or, when $B<P$, the tangent kernel
$K=OO^{\!\top}$, whose nonzero spectrum equals that of $S$ up to the overall factor $B$). From the spectrum
$\{\lambda_a\}$ we report the effective dimension
$d_{\mathrm{eff}}=(\sum_a\lambda_a)^2/\sum_a\lambda_a^2$ and the condition number
$\kappa(S)=\lambda_{\max}/\lambda_{\min}$. Cross-architecture comparisons of
$d_{\mathrm{eff}}$ are made on a \emph{common probe set} (all models scored on the same configurations, not on each model's own $|\Psi|^2$) so that the comparison reflects the
ansatz and not its sampling measure, and we report only participation ratios that have
converged in the batch size.

\paragraph{A physical-operator dictionary.}
Counting tangent directions says how many; naming them says which. We therefore project the
leading tangent modes onto a small, fixed dictionary of one- and two-body operators chosen
before the fits were run, each written as a scalar function of a configuration whose
parameter-space image is obtained by the same score construction as $O$:
\begin{equation}
\begin{aligned}
  \mathcal{D}_1 &= \textstyle\sum_i \tilde r_i^{\,2}
    &&\text{breathing / monopole, the generator of $\partial_\omega$;}\\
  \mathcal{D}_2 &= \textstyle\sum_i \tilde r_i^{\,2}\cos 2\theta_i
    &&\text{quadrupole, the lowest shape distortion;}\\
  \mathcal{D}_3 &= \textstyle\sum_{i<j} e^{-\tilde r_{ij}^{2}/2\sigma_h^{2}}
    &&\text{correlation hole, with $\sigma_h$ the fitted hole width;}\\
  \mathcal{D}_4 &= \textstyle\sum_{i<j} \tilde r_{ij}
    &&\text{pair-linear, the leading long-range pair response.}
\end{aligned}
\label{eq:dictionary}
\end{equation}
Writing $D\in\mathbb{R}^{B\times4}$ for the centred matrix of these four functions on the
probe batch and $v_a\in\mathbb{R}^{B}$ for the $a$th tangent-kernel eigenvector, the
\emph{naming} score of a mode is the coefficient of determination of its least-squares
regression on the dictionary,
\begin{equation}
  R^2_a \;=\; 1-\frac{\|(I-P_D)\,v_a\|^2}{\|v_a\|^2},
  \qquad P_D=D(D^{\!\top}D)^{-1}D^{\!\top},
  \label{eq:mode-naming}
\end{equation}
and the \emph{dictionary-orthogonal tangent mass} is the eigenvalue-weighted share of the
leading $A$ modes lying outside that span,
\begin{equation}
  M_\perp \;=\; \frac{\sum_{a\le A}\lambda_a\,\|(I-P_D)v_a\|^2}{\sum_{a\le A}\lambda_a},
  \qquad A=10 .
  \label{eq:nonphys-mass}
\end{equation}
Finally, \emph{correlation-hole capture} is $R^2$ of Eq.~\eqref{eq:mode-naming} computed
against the single operator $\mathcal{D}_3$, i.e.\ the share of the leading mode explained
by the hole channel alone.

The dictionary is a \emph{chosen, non-exhaustive} probe, and nothing here establishes that
it spans the physically meaningful directions. A large $M_\perp$ means the tangent space is
poorly accounted for by these four operators; it does not mean the remaining directions are
unphysical. The comparison supports a relative statement between two architectures
scored against the same dictionary on the same probe batch, and that is how
Chapter~\ref{ch:kernel} uses it.

\paragraph{Reporting the condition number.}
$\kappa(S)=\lambda_{\max}/\lambda_{\min}$ is quoted with two caveats that its definition
forces. Our batches have $B<P$, so $S$ is rank-deficient by construction and $\lambda_{\min}$
is the smallest \emph{sample-supported} eigenvalue, not an intrinsic one; the value of
$\kappa$ therefore depends on $B$ and is a property of the empirical metric the optimiser
actually sees, not of the variational manifold. That is the relevant object for the question we ask of it, whether preconditioning by $S^{-1}$ helps, but it means $\kappa$ may not be
compared across batch sizes, and we hold $B$ fixed wherever we compare it.

To ask whether two trained models are the \emph{same} state rather than merely equal in
energy, we use a symmetric, importance-sampled squared overlap,
\[
  \mathcal{S}^2(A,B)
  = \big\langle \Psi_B/\Psi_A\big\rangle_{|\Psi_A|^2}\;
    \big\langle \Psi_A/\Psi_B\big\rangle_{|\Psi_B|^2},
\]
each factor evaluated by a log-sum-exp mean over samples of the corresponding state.
It equals $1$ for identical states, and on the analytic two-electron ground state the
trained ansatz reproduces $\mathcal{S}^2=1.000000$, which anchors the tooling.

One convention must be stated, because the ratio $\Psi_B/\Psi_A$ is signed for a
fermionic state. We evaluate the ratios from $\log|\Psi|$, so the estimator is strictly a
\emph{modulus} overlap, $\langle|\Psi_B/\Psi_A|\rangle\langle|\Psi_A/\Psi_B|\rangle$. For
the comparisons made here this is an accurate proxy for the signed overlap: both states
share the same Slater reference, and the correlator and cusp enter as a strictly positive
exponential, so the sign of $\Psi$ is fixed by the determinant and the two states agree in
sign except on the measure-suppressed set where their backflows carry a configuration
across a node differently. The modulus overlap is an \emph{upper} bound on the signed one, so $\mathcal S^2\to0$ is
conservative while $\mathcal S^2\approx1$ rests on the structural argument above rather than
on the bound. Overlap, energy and $\Var(E_L)$ are properties of the state and are unaffected
by how correlation is divided between the correlator and the backflow; the energetic price of
ablating a component is not, since two parameterisations that divide correlation differently
report different ablation costs.
